% =====================================================================
%  Requirements Specification -- Recommender-enabled e-commerce platform
%  Stakeholders, Functional Requirements and Use-Case Model
%  Course: Intelligent Systems (CO5119) -- HCMUT
%  Compile: pdflatex main.tex  (run twice for TOC / cross-references)
% =====================================================================
\documentclass[a4paper,12pt]{article}

% ---------- Encoding and fonts ----------
% Works with pdfLaTeX (T5 encoding, vntex fonts, for the Vietnamese name on the
% title page) and with XeLaTeX / LuaLaTeX (fontspec, Latin Modern).
\usepackage{iftex}
\ifPDFTeX
  \usepackage[utf8]{inputenc}
  \usepackage[T5]{fontenc}
\else
  \usepackage{fontspec}   % defaults to Latin Modern, which covers Vietnamese glyphs
\fi
\usepackage[english]{babel}

% ---------- Page layout ----------
\usepackage[a4paper,top=2.5cm,bottom=2.5cm,left=2.5cm,right=2.5cm]{geometry}
\usepackage{setspace}
\onehalfspacing
\usepackage{parskip}

% ---------- Graphics, tables, lists ----------
\usepackage{graphicx}
\graphicspath{{./}{figures/}}
\usepackage{longtable}
\usepackage{array}
\usepackage{tabularx}
\usepackage{ragged2e}
\usepackage{enumitem}
\usepackage[table]{xcolor}
\usepackage{caption}
\usepackage{float}
\captionsetup{font=small,labelfont=bf}

% ---------- Header / footer ----------
\usepackage{fancyhdr}
\pagestyle{fancy}
\fancyhf{}
\fancyhead[L]{\small Requirements Specification}
\fancyhead[R]{\small Intelligent Systems (CO5119)}
\fancyfoot[C]{\small Page \thepage\ of \pageref{LastPage}}
\setlength{\headheight}{16pt}
\renewcommand{\headrulewidth}{0.4pt}
\usepackage{lastpage}

% ---------- Hyperlinks ----------
\usepackage[hidelinks]{hyperref}
\hypersetup{
  pdftitle={Requirements Specification -- Recommender-enabled e-commerce platform},
  pdfsubject={Stakeholders, Functional Requirements, Use-Case Model},
  pdfkeywords={association rule mining, word embedding, recommender system, requirements}
}

% ---------- Helpers ----------
\definecolor{accent}{HTML}{E9562A}
\definecolor{headbg}{HTML}{FFE9D6}
\definecolor{tableborder}{HTML}{808080}
\newcolumntype{L}[1]{>{\RaggedRight\arraybackslash}p{#1}}
\newcolumntype{C}[1]{>{\centering\arraybackslash}p{#1}}
\newlength{\tablecontentwidth}
\setlength{\tablecontentwidth}{\dimexpr\textwidth-8\tabcolsep-5\arrayrulewidth\relax}
\newcommand{\tablestyle}{%
  \small\setstretch{1.05}%
  \renewcommand{\arraystretch}{1.25}%
  \setlength{\extrarowheight}{3pt}%
  \arrayrulecolor{tableborder}%
}
\setlength{\LTpre}{0.6\baselineskip}
\setlength{\LTpost}{0.6\baselineskip}
\newcommand{\code}[1]{\texttt{#1}}
\setlength{\emergencystretch}{3em}   % avoid overfull lines around long code identifiers
\newcommand{\req}[1]{\textsf{\small #1}}
\newcommand{\inc}{\textit{$\ll$include$\gg$}}

% Preserve the original assets in project builds. Standalone previews that cannot
% read companion files display an explicit placeholder instead of failing.
\newcommand{\reportimage}[4][]{%
  \IfFileExists{#2}{%
    \includegraphics[#1]{#2}%
  }{%
    \PackageWarningNoLine{report}{Image #2 unavailable in this compiler}%
    \begingroup
    \setlength{\fboxsep}{6pt}%
    \fbox{\parbox[c][#4][c]{\dimexpr#3-2\fboxsep-2\fboxrule\relax}{%
      \centering\small\normalfont
      \textbf{Image unavailable in this preview}\\[0.6em]
      \texttt{\detokenize{#2}}\\[0.6em]
      Compile with the accompanying image files to display the original.%
    }}%
    \endgroup
  }%
}

% =====================================================================
\begin{document}
\hypersetup{pageanchor=false}

% ---------------------------------------------------------------------
%  Title page
% ---------------------------------------------------------------------
\begin{titlepage}
  \centering
  \setstretch{1.1}
  {\normalsize\bfseries VIETNAM NATIONAL UNIVERSITY -- HO CHI MINH CITY\\
   HO CHI MINH CITY UNIVERSITY OF TECHNOLOGY\\
   FACULTY OF COMPUTER SCIENCE AND ENGINEERING\par}
  \vspace{0.4cm}
  \reportimage[width=0.24\textwidth]{hcmut.png}{0.24\textwidth}{0.24\textwidth}
  \par\vspace{0.4cm}

  {\LARGE\bfseries REQUIREMENTS SPECIFICATION\par}
  \vspace{0.6cm}
  {\large\bfseries SUBJECT: INTELLIGENT SYSTEMS (CO5119)\par}
  \vspace{0.6cm}
  \begin{minipage}{0.94\textwidth}
    \centering\large\bfseries
    TOPIC: A HYBRID PRODUCT RECOMMENDER FOR AN ELECTRONICS E-COMMERCE
    PLATFORM USING ASSOCIATION RULE MINING AND WORD EMBEDDINGS\par
  \end{minipage}\par
  \vspace{1.2cm}

  {\large Supervisor: \textbf{Assoc.\ Prof.\ Dr.\ Quản Thành Thơ}\par}
  \vspace{0.6cm}
  {\large Students:\par}
  \vspace{0.2cm}
  {\tablestyle
  \begin{tabularx}{\textwidth}{|L{3.6cm}|C{1.6cm}|>{\RaggedRight\arraybackslash}X|}
    \hline
    \rowcolor{headbg}
    \textbf{Name} & \textbf{ID} & \textbf{Email} \\[3pt]\hline
    Phan Hữu Thành & 2670371 & thanh.phanbluewhales@hcmut.edu.vn \\[3pt]\hline
    Trần Anh Khoa & 2670309 & khoa.trana@hcmut.edu.vn \\[3pt]\hline
  \end{tabularx}
  }

  \vfill
  {\large Ho Chi Minh City, 10/2026\par}
\end{titlepage}
\pagenumbering{roman}
\hypersetup{pageanchor=true}

% ---------------------------------------------------------------------
%  Front matter
% ---------------------------------------------------------------------
\tableofcontents
\clearpage
\listoffigures
\listoftables
\clearpage
\pagenumbering{arabic}

% ---------------------------------------------------------------------
%  Scope
% ---------------------------------------------------------------------
\section*{Scope of this document}
\addcontentsline{toc}{section}{Scope of this document}

This document specifies the requirements of a recommender-enabled electronic e-commerce platform developed as
the laboratory platform for the course topics \emph{Association Rules Mining}, \emph{Recommender
Systems} and \emph{Word Embedding}. The e-commerce platform, the administration dashboard and the data layer
constitute the operational part of the system. The intelligent component, a hybrid recommender that
combines association rules learned from purchase transactions with semantic similarity learned from
product descriptions, constitutes the research part, together with the data-warehouse pipeline that
prepares the training data on which the recommender depends.

The specification answers three questions: \emph{who} interacts with the system, \emph{what} the
system must do for them, and \emph{how} these capabilities are organised into use cases.
Section~\ref{sec:stakeholders} describes the role and benefits of each stakeholder group.
Section~\ref{sec:fr} lists the functional requirements with their priority, including the
requirements of the data-warehouse pipeline that supplies training data to the models.
Section~\ref{sec:usecase} presents the general use-case model, the use-case catalogue and the main
flows of the central use cases. Section~\ref{sec:traceability} provides the requirement-to-use-case
traceability mapping in Table~\ref{tab:traceability}, covering both requirements and catalogue
entries. Section~\ref{sec:dataset} introduces the selected dataset,
describes its files and fields, and presents the adapted high-level operational data model.
Section~\ref{sec:prototype} presents the interactive prototype and interface mockups,
distinguishing demonstrated user interactions from the remaining analytical implementation.
Detailed warehouse schemas,
programming interfaces and deployment procedures are documented separately in the project's
architecture description. Warehouse-schema design and ETL authoring are project responsibilities
carried out through external scripts and notebooks; the functional requirements describe execution
of the resulting pipeline and the recommendation workflow supported by the system.

For this project, the model-development data consists of completed purchase transactions and
available product text. Orders retain their items, quantities, purchase-time prices, timestamps and
status; the catalogue retains product identifiers, categories and available names, descriptions and
tags. Versioned datasets, training parameters, model artefacts and offline evaluation results make
experiments repeatable. The current cart supplies recommendation context, but historical product
views, clicks, cart actions and recommendation impressions are outside the recording scope.
Purchases are retained as transactions without attribution to recommendations. Dataset refresh and
model improvement are initiated manually by the data scientist.

% ---------------------------------------------------------------------
\clearpage
\section{Role and Benefits of Stakeholders}
\label{sec:stakeholders}

The application applies Association Rule Mining to purchase transactions to discover products that
are frequently bought together, and Word Embeddings to product text to identify semantically related
products. Three stakeholder groups interact with it: \textbf{product buyers}, \textbf{administrators}
and \textbf{data scientists}. Buyers generate purchase transactions, data scientists consolidate
them with catalogue data in a warehouse and turn them into models, and administrators operate the
store, safeguard the input data and inspect the active model. Table~\ref{tab:stakeholders}
summarises the three groups.

\begin{table}[htbp]
  \centering
  \caption{Summary of stakeholders, their touch-points and primary benefits.}
  \label{tab:stakeholders}
  \tablestyle
  \begin{tabular}{|L{0.16\tablecontentwidth}|L{0.28\tablecontentwidth}|L{0.22\tablecontentwidth}|L{0.34\tablecontentwidth}|}
    \hline
  \rowcolor{headbg}
    \textbf{Stakeholder} & \textbf{Main role} & \textbf{Touch-point} & \textbf{Primary benefit} \\[3pt]\hline
    Product Buyer & Browses, compares and purchases products; generates purchase transactions &
      Public storefront & Finds suitable products faster; receives explainable
      ``bought together'' and ``similar product'' suggestions \\[3pt]\hline
    Administrator & Manages the catalogue, views orders, tracks sales, safeguards data quality, views the
      active model & Administration dashboard & Operates the store from a single interface;
      automated product suggestions; visibility of the active model \\[3pt]\hline
    Data Scientist & Runs the warehouse pipeline, mines rules, trains embeddings, evaluates and
      publishes models, refreshes experiments and selects the active configuration & Analytics workspace (external notebooks
      and scripts, data
      warehouse) & A realistic case study with clean, versioned data and a measurable evaluation
      loop \\[3pt]\hline
  \end{tabular}
\end{table}

\subsection{Product Buyers}

Product buyers are the end consumers of the storefront: they browse categories, search by keyword,
inspect product details, maintain a cart and complete purchases, without registration. Their role
is to select products and respond to the recommendations shown to them. Opening a product page
triggers rule-based ``frequently bought together'' suggestions and embedding-based ``similar
product'' suggestions; changing the cart recomputes them over the whole basket. Completed purchases
are recorded as orders and line items with quantities, purchase-time prices,
timestamps and status. These transactions, together with maintained catalogue text, provide data for
offline training and evaluation alongside the imported historical purchases. Cart contents are
persisted as current shopping state and used to request suggestions; they are not an interaction
history. Product views and recommendation displays or clicks are not recorded for model development.

Buyers benefit from data-driven suggestions instead of scanning the whole catalogue: association
rules reveal complementary items that are easy to overlook, embeddings reveal substitutes when a
product is unavailable or too expensive, and each suggestion carries a short explanation that
builds trust. A persistent cart and a one-step checkout reduce the effort of completing an order.

\subsection{Administrators}

Administrators operate and maintain the store through the administration dashboard; they understand
the retail business and need no machine-learning background. Their responsibilities are catalogue
management, including withdrawing products without losing order history; viewing orders and their
contained products; monitoring sales through
revenue, order volume, average order value and best-sellers; and, in the extended system, overseeing
the recommender in production by viewing which model version or configuration is active. Because
the embedding model learns from available product names, descriptions, categories and tags,
administrators are also indirectly responsible for the quality of that text.

Administrators benefit from a single, validated interface for all operations, from reports that
remain correct when products are withdrawn, and from automated product suggestions without manually
curating product pairs. Serving reports from the data warehouse keeps long analytical queries away from the
storefront, and the modular design makes the system easier to maintain and extend.

\subsection{Data Scientists}

Data scientists, including the students carrying out the laboratory, are responsible for the
analytical and experimental side of the application. As project responsibilities performed through
external scripts and notebooks, they \emph{design the data
warehouse schema}: identify the purchase processes to analyse, fix the grain of each fact table
(one purchased order line), define its measures and the conformed product, category and calendar-date
dimensions, including surrogate keys and the handling of changing product attributes such as price
and category, and author and modify the ETL pipeline. These development activities are outside the
application's use-case boundary. Through the analytics workspace they execute and refresh the
resulting pipeline, and extract versioned datasets from successful warehouse loads: distinct-product
baskets from completed orders for rule mining, an available product-text corpus for embeddings, and
chronologically separated purchases for offline evaluation.

On these datasets they mine association rules with Apriori or FP-Growth, tuning support, confidence
and lift; train word embeddings with Word2Vec or FastText, or use a sentence-embedding model, to obtain
one vector per product; and configure how association-based and embedding-based candidates are
combined, with popularity as a fallback when learned candidates are insufficient. Rule-only,
embedding-only and hybrid approaches are evaluated separately using suitable offline evidence.
Purchase-basket completion evaluates rule-based and hybrid recommendations against a popularity
baseline using precision, recall and catalogue coverage at a fixed list length on later held-out
orders. Similar-product embeddings are assessed through a documented manual relevance review of a
fixed sample against a category-based baseline, since substitutes need not be bought together.
The data scientist stores the versioned artefacts and evaluation results, publishes versions that meet the
acceptance criteria, and has sole responsibility for selecting and activating the model or
configuration used by the storefront. Publishing makes an evaluated version available; selection
determines which available version is served. Both actions may be performed consecutively by the
same data scientist. When new completed purchases or catalogue changes are available, the data
scientist manually refreshes the warehouse and datasets, reruns training and evaluation, and compares
the new results before deciding whether to publish and activate a replacement.

Data scientists benefit from a realistic e-commerce case study in which transactional and textual
data come from one live system, from a warehouse whose schema is shaped in advance for rule mining,
embedding and evaluation, and from snapshot-based training that isolates experiments from operations
and makes them repeatable and comparable. Stored dataset versions, training parameters, model
versions and evaluation results support model improvement without an online interaction-feedback
pipeline or claims about purchases caused by recommendations.

% ---------------------------------------------------------------------
\clearpage
\section{Functional Requirements}
\label{sec:fr}

The following tables group functional requirements by stakeholder and pipeline.
Each requirement has an identifier: FR-B for buyers, FR-A for administrators, FR-D for data scientists
and FR-W for the warehouse pipeline. Table~\ref{tab:traceability} maps these identifiers to the
use-case catalogue. Schema design and ETL authoring remain the project responsibilities described
in Section~\ref{sec:stakeholders}, rather than application functional requirements.
Priority follows the MoSCoW scheme: \textbf{M}ust have, \textbf{S}hould have and
\textbf{C}ould have.
    
\subsection{Requirements for Product Buyers}

{\tablestyle
\begin{longtable}{|L{\dimexpr\textwidth-2cm-4\tabcolsep-3\arrayrulewidth\relax}|C{2cm}|}
  \caption{Functional requirements for product buyers.}
  \label{tab:fr-buyer}\\
  \hline
  \rowcolor{headbg}
  \textbf{Requirement} & \textbf{Priority} \\
  \hline
  \endfirsthead
  \hline
  \rowcolor{headbg}
  \textbf{Requirement} & \textbf{Priority} \\
  \hline
  \endhead
  \endfoot
  \endlastfoot

  \req{FR-B01} The buyer can browse, search and view products available on the e-commerce platform. & M \\[3pt]\hline
  \req{FR-B02} The buyer can add products to a shopping cart, update the cart and complete a purchase. & M \\[3pt]\hline
  \req{FR-B03} Completed purchases are stored as orders and line items retaining order and product
  identifiers, quantities, purchase-time prices, purchase timestamps and status. These records support
  purchase-basket mining, offline evaluation and sales reporting without recommendation attribution. & M \\[3pt]\hline
  \req{FR-B04} The buyer receives ``frequently bought together'' recommendations derived from association
  patterns learned from historical completed purchase transactions. & M \\[3pt]\hline
  \req{FR-B05} The buyer receives additional ``similar product'' recommendations derived from semantic
  similarity between product representations produced using word embeddings. & M \\[3pt]\hline
  \req{FR-B06} Recommendations are presented with a short explanation indicating why a product is suggested,
  such as being frequently associated with products in the current context or being semantically
  similar to the product being viewed. & S \\[3pt]\hline
\end{longtable}
}


\subsection{Requirements for Administrators}

{\tablestyle
\begin{longtable}{|L{\dimexpr\textwidth-2cm-4\tabcolsep-3\arrayrulewidth\relax}|C{2cm}|}
  \caption{Functional requirements for administrators.}
  \label{tab:fr-admin}\\
  \hline
  \rowcolor{headbg}
  \textbf{Requirement} & \textbf{Priority} \\
  \hline
  \endfirsthead
  \hline
  \rowcolor{headbg}
  \textbf{Requirement} & \textbf{Priority} \\
  \hline
  \endhead
  \endfoot
  \endlastfoot

  \req{FR-A01} The administrator can create, update and withdraw products, including their catalogue
  information such as name, category, price, stock level, description and tags. These operational
  changes are stored in the operational database. & M \\[3pt]\hline
  \req{FR-A02} The administrator can view orders and the products contained in each order. & M \\[3pt]\hline
  \req{FR-A03} The administrator can view analytical information derived from the data warehouse, including
  revenue, order volume, average order value, best-selling products and sales trends over time. & M \\[3pt]\hline
  \req{FR-A04} The administrator can view which recommendation model or recommendation configuration is
  currently active. & S \\[3pt]\hline
\end{longtable}
}


\subsection{Requirements for Data Scientists}

{\tablestyle
\begin{longtable}{|L{\dimexpr\textwidth-2cm-4\tabcolsep-3\arrayrulewidth\relax}|C{2cm}|}
  \caption{Functional requirements for data scientists.}
  \label{tab:fr-ds}\\
  \hline
  \rowcolor{headbg}
  \textbf{Requirement} & \textbf{Priority} \\
  \hline
  \endfirsthead
  \hline
  \rowcolor{headbg}
  \textbf{Requirement} & \textbf{Priority} \\
  \hline
  \endhead
  \endfoot
  \endlastfoot

  \req{FR-D01} The data scientist can manually execute and refresh the ETL pipeline that transforms
  operational product and purchase-transaction data into the analytical structures stored in the
  data warehouse. & M \\[3pt]\hline
  \req{FR-D02} The data scientist can obtain versioned completed-order baskets, available product-text
  corpora and chronological purchase hold-outs from the data warehouse for mining, semantic product
  representation and offline recommendation experiments. & M \\[3pt]\hline
  \req{FR-D03} The data scientist can apply Association Rule Mining to historical completed purchase
  baskets to discover product associations used for ``frequently bought together'' recommendations. & M \\[3pt]\hline
  \req{FR-D04} The data scientist can use word-embedding techniques to produce semantic representations of
  products from their textual information and identify semantically similar products. & M \\[3pt]\hline
  \req{FR-D05} The data scientist can combine association-rule-based candidates and embedding-based candidates
  into a hybrid recommendation strategy, while also retaining the ability to analyse the two
  recommendation approaches separately. & M \\[3pt]\hline
  \req{FR-D06} The data scientist can store versioned association rules, product representations, hybrid
  configuration and model metadata, including source dataset versions and training parameters, so
  that experiments can be reproduced and the artefacts evaluated, published and used by the platform. & M \\[3pt]\hline
  \req{FR-D07} The data scientist can evaluate rule-based and hybrid basket-completion recommendations
  on later held-out purchases using precision, recall and catalogue coverage at a fixed list length
  against a popularity baseline. Embedding-based similar products are assessed through a documented
  manual relevance review of a fixed sample against a category-based baseline. Metrics and review
  results are stored with the corresponding model and dataset versions. The evaluation protocol and
  acceptance criteria are recorded before evaluating a candidate. & M \\[3pt]\hline
  \req{FR-D08} The data scientist can publish evaluated model versions that meet the defined acceptance criteria,
  releasing their association rules, product representations, hybrid configuration and model metadata
  so that they are available for selection and activation by the data scientist. & M \\[3pt]\hline
  \req{FR-D09} The data scientist can select and activate an available published recommendation model or
  configuration for use by the e-commerce platform, and view confirmation of the active selection.
  Model selection and activation are restricted to this role. & S \\[3pt]\hline
  \req{FR-D10} After new completed purchases or catalogue updates, the data scientist can manually
  refresh the datasets, retrain affected models and rerun offline evaluation. The resulting versions
  and results remain available for comparison before publication and activation. & S \\[3pt]\hline
\end{longtable}
}


\subsection{Requirements for the Data-Warehouse Pipeline}

The data warehouse supports both the intelligent recommendation workflow and the analytical
functions of the administration dashboard. It consolidates operational data into a structure that
can be used for Association Rule Mining, word-embedding-based recommendation and business
statistics.

{\tablestyle
\begin{longtable}{|L{\dimexpr\textwidth-2cm-4\tabcolsep-3\arrayrulewidth\relax}|C{2cm}|}
  \caption{Functional requirements for the data-warehouse pipeline.}
  \label{tab:fr-dw}\\
  \hline
  \rowcolor{headbg}
  \textbf{Requirement} & \textbf{Priority} \\
  \hline
  \endfirsthead
  \hline
  \rowcolor{headbg}
  \textbf{Requirement} & \textbf{Priority} \\
  \hline
  \endhead
  \endfoot
  \endlastfoot

  \req{FR-W01} The ETL pipeline extracts relevant product, order and purchased-item data from the
  operational database and loads the transformed data into a separate analytical data warehouse. & M \\[3pt]\hline
  \req{FR-W02} The data warehouse contains fact and dimension structures that preserve purchased
  items, quantities, purchase-time prices, order status, product information and purchase timestamps
  required by the recommender system and administrative reports. & M \\[3pt]\hline
  \req{FR-W03} The warehouse provides distinct-product baskets grouped by eligible completed order
  for Association Rule Mining. Quantities remain available for sales reporting; repeated units of
  the same product do not create duplicate basket entries. & M \\[3pt]\hline
  \req{FR-W04} The warehouse provides available product text, including categories and any maintained
  names, descriptions and tags, for embedding-based representations. The initial Olist dataset uses
  translated category text; any richer project catalogue text is documented separately. & M \\[3pt]\hline
  \req{FR-W05} The warehouse provides aggregated historical information required by the administration
  dashboard, including sales, orders and product performance over time. & M \\[3pt]\hline
  \req{FR-W06} The data scientist can manually rerun the pipeline after new purchase transactions or
  catalogue changes to refresh the warehouse and the datasets used by the recommendation system. & S \\[3pt]\hline
\end{longtable}
}
% ---------------------------------------------------------------------
\clearpage
\section{General Use-Case Model}
\label{sec:usecase}

Figure~\ref{fig:usecase} shows the general use-case model. The three actors operate in three
subsystems that are connected through data rather than through direct interaction: buyers create
purchase transactions, the warehouse pipeline consolidates them with catalogue data, data scientists extract training
datasets from the warehouse, and the model selected and activated by the data scientist supplies the
recommendations shown to buyers. Administrators view orders and sales reports, maintain the catalogue
text on which embeddings depend, and view the active model or configuration.
The analytics workspace exposes pipeline execution and model operations; warehouse-schema design
and ETL authoring take place in external development tools.

Solid lines associate actors with their use cases. Grey dashed arrows labelled \inc\ denote
behaviour performed within another use case, namely recommendation delivery on product pages and
cart updates. Purple dashed arrows denote data dependencies across subsystem boundaries.
They describe the flow of model, purchase-transaction and catalogue information, rather than
the order in which use cases must execute. Previously prepared warehouse snapshots, datasets and
evaluation results may be reused; their availability is a precondition for subsequent operations.

\begin{figure}[htbp]
  \centering
  \reportimage[width=\textwidth,height=0.82\textheight,keepaspectratio]{figures/usecase.png}{\textwidth}{0.81\textwidth}
  \caption{General use-case diagram of the recommender-enabled storefront.}
  \label{fig:usecase}
\end{figure}
\clearpage

\subsection{Use-case catalogue}

Table~\ref{tab:usecases} summarises each use case with its primary actor and outcome.

{\tablestyle
\begin{longtable}{|L{0.20\tablecontentwidth}|L{0.14\tablecontentwidth}|L{0.33\tablecontentwidth}|L{0.33\tablecontentwidth}|}
  \caption{Use-case catalogue.}
  \label{tab:usecases}\\
  \hline
  \rowcolor{headbg}
  \textbf{Use case} & \textbf{Actor} & \textbf{Summary} & \textbf{Outcome} \\
  \hline
  \endfirsthead
  \hline
  \rowcolor{headbg}
  \textbf{Use case} & \textbf{Actor} & \textbf{Summary} & \textbf{Outcome} \\
  \hline
  \endhead
  \endfoot
  \endlastfoot
  Browse and search products & Buyer & List products by category; search, filter and sort &
    Matching products displayed \\[3pt]\hline
  View product details & Buyer & Open a product page; includes \emph{Receive recommendations} &
    Product attributes and recommendations displayed \\[3pt]\hline
  Receive recommendations & Buyer (indirect) & Present complementary and similar suggestions for
    the current product or basket & Explained list generated using the active model and strategy \\[3pt]\hline
  Manage shopping cart & Buyer & Add, update and remove items; includes \emph{Receive
    recommendations} & Current cart updated and persisted; recommendations refreshed \\[3pt]\hline
  Place an order & Buyer & Confirm the cart with minimal contact details & Order and line items
    recorded atomically; stock reduced; cart cleared \\[3pt]\hline

  View orders and order items & Administrator & Inspect an order and its contained products &
    Order details and historical item information displayed \\[3pt]\hline
  View sales reports & Administrator & Inspect warehouse-derived revenue, order volume, average
    order value, best-sellers and trends & Requested sales indicators and time series displayed \\[3pt]\hline
  Manage catalogue & Administrator & Create, edit and withdraw products & Catalogue updated;
    withdrawn products kept in historical records \\[3pt]\hline
  View model status & Administrator & Inspect the active recommendation model or configuration &
    Active version and configuration displayed \\[3pt]\hline

  Run and refresh data warehouse & Data Scientist & Execute ETL into the established dimensional
    schema & Facts, dimensions, dataset views and reporting aggregates refreshed; load metadata
    and quality report recorded \\[3pt]\hline
  Extract dataset & Data Scientist & Derive mining, product-text and evaluation data from a
    successful warehouse snapshot & Versioned datasets stored with lineage \\[3pt]\hline
  Train models & Data Scientist & Mine rules, train embeddings and configure the hybrid using
    available versioned datasets & Versioned rules, vectors, hybrid configuration and metadata
    stored; separate methods remain available \\[3pt]\hline
  Evaluate models & Data Scientist & Assess basket completion on held-out purchases and review
    similar-product relevance against suitable baselines & Offline metrics and manual review results
    stored with model and dataset versions \\[3pt]\hline
  Publish evaluated model & Data Scientist & Release stored rules, vectors, hybrid configuration
    and metadata after checking acceptable evaluation results & Published version available for
    selection and activation \\[3pt]\hline
  Select and activate model/configuration & Data Scientist & Select an available published version
    or configuration for online serving & Active selection updated and confirmed; subsequent
    recommendations use it \\[3pt]\hline
\end{longtable}
}

\subsection{Requirement-to-use-case traceability}
\label{sec:traceability}

Table~\ref{tab:traceability} links every functional requirement to the use case or use cases that
describe its behaviour. A use case may cover several requirements, and a requirement may involve
several use cases. Purchase persistence is part of placing an order; warehouse refresh and model
experiments are separate Data Scientist operations. Schema design and ETL authoring are outside this
mapping because they are project responsibilities rather than application functional requirements.
Every catalogue entry appears in the mapping, so it can also be read from a use case to the
requirements it fulfils.

{\tablestyle
\begin{longtable}{|L{2.6cm}|L{\dimexpr\textwidth-2.6cm-4\tabcolsep-3\arrayrulewidth\relax}|}
  \caption{Functional requirement-to-use-case traceability.}
  \label{tab:traceability}\\
  \hline
  \rowcolor{headbg}
  \textbf{Requirement ID} & \textbf{Corresponding use case(s)} \\
  \hline
  \endfirsthead
  \hline
  \rowcolor{headbg}
  \textbf{Requirement ID} & \textbf{Corresponding use case(s)} \\
  \hline
  \endhead
  \endfoot
  \endlastfoot
  \req{FR-B01} & Browse and search products; View product details \\[3pt]\hline
  \req{FR-B02} & Manage shopping cart; Place an order \\[3pt]\hline
  \req{FR-B03} & Place an order \\[3pt]\hline
  \req{FR-B04} & Receive recommendations \\[3pt]\hline
  \req{FR-B05} & Receive recommendations \\[3pt]\hline
  \req{FR-B06} & Receive recommendations \\[3pt]\hline
  \req{FR-A01} & Manage catalogue \\[3pt]\hline
  \req{FR-A02} & View orders and order items \\[3pt]\hline
  \req{FR-A03} & View sales reports \\[3pt]\hline
  \req{FR-A04} & View model status \\[3pt]\hline
  \req{FR-D01} & Run and refresh data warehouse \\[3pt]\hline
  \req{FR-D02} & Extract dataset \\[3pt]\hline
  \req{FR-D03} & Train models \\[3pt]\hline
  \req{FR-D04} & Train models \\[3pt]\hline
  \req{FR-D05} & Train models; Evaluate models \\[3pt]\hline
  \req{FR-D06} & Train models; Publish evaluated model \\[3pt]\hline
  \req{FR-D07} & Evaluate models \\[3pt]\hline
  \req{FR-D08} & Publish evaluated model \\[3pt]\hline
  \req{FR-D09} & Select and activate model/configuration \\[3pt]\hline
  \req{FR-D10} & Run and refresh data warehouse; Extract dataset; Train models; Evaluate models \\[3pt]\hline
  \req{FR-W01} & Run and refresh data warehouse \\[3pt]\hline
  \req{FR-W02} & Run and refresh data warehouse \\[3pt]\hline
  \req{FR-W03} & Run and refresh data warehouse; Extract dataset \\[3pt]\hline
  \req{FR-W04} & Run and refresh data warehouse; Extract dataset \\[3pt]\hline
  \req{FR-W05} & Run and refresh data warehouse; View sales reports \\[3pt]\hline
  \req{FR-W06} & Run and refresh data warehouse; Extract dataset \\[3pt]\hline
\end{longtable}
}

\subsection{Main flows of the central use cases}

\paragraph{Receive recommendations.}
The storefront requests suggestions for the current product or cart contents. Using the active model
and configuration, the recommendation service retrieves complementary candidates from matching
association rules and similar-product candidates from product embeddings. The hybrid configuration
determines their ranking and combination; popular products fill remaining gaps when learned
candidates are insufficient. Each suggestion carries its model version, strategy and a short reason.
The storefront displays the suggestions with their reasons. Product-page and cart changes request
updated suggestions using the current context; they do not create historical interaction records.

\paragraph{Place an order.}
The buyer confirms the cart with minimal contact details. The order service validates the request and
checks availability, then records the order, its line items with the price at the time of purchase,
and the stock adjustments in one atomic operation. The order retains its purchase timestamp and
status, and the buyer receives an order number. Completed orders become the durable transaction
history used by the warehouse; no purchase-to-recommendation attribution is recorded. If any item is
insufficiently stocked, the whole operation is cancelled and the
affected product is reported.

\paragraph{View orders and order items.}
The administrator selects an order from the operational order records. The system displays its
details and contained products, including quantities and prices recorded at purchase time.
Historical order information remains accessible when a catalogue product is withdrawn.

\paragraph{View sales reports.}
The administrator chooses the reporting period and inspects warehouse-derived revenue, order volume,
average order value, best-selling products and sales trends. Reports reflect the latest successful
warehouse refresh and retained purchase-time quantities and prices.

\paragraph{View model status.}
The administrator inspects the currently active model version and recommendation configuration from
the serving state. This visibility is separate from sales reporting and does not grant model
selection or activation rights.

\paragraph{Run and refresh data warehouse.}
When manually requested by the data scientist, the pipeline extracts orders, purchased items and
products changed since the last successful load. It transforms them to the dimensional schema
(surrogate keys, conformed dimensions and facts at purchased-order-line grain), validates integrity,
completeness, duplicates and value ranges, and loads the result into the schema prepared during
project development. Failed checks reject the load; successful loads refresh training-dataset views
and sales aggregates used by the administration dashboard, and record their lineage. Subsequent
runs process new purchase transactions and catalogue changes. Executing ETL is separate from
designing or modifying its implementation.

\paragraph{Extract dataset.}
A successful warehouse snapshot is a precondition. The data scientist extracts distinct-product
baskets from eligible completed orders, available product text including categories and any
maintained names, descriptions and tags, and chronological purchase data for evaluation. Imported
Olist records initially supply purchase baskets and translated category text; richer product text,
if added, comes from separately documented catalogue enrichment. Quantities are obtained from stored
order-item quantities or counted source item rows as appropriate. Each dataset is stored with its
source snapshot, scope and version. Extraction can reuse an existing warehouse snapshot without
rerunning ETL.

\paragraph{Model life cycle.}
Training requires available versioned datasets. The data scientist mines rules with chosen support,
confidence and lift thresholds, produces one embedding vector per product, and configures the hybrid
combination and fallback behaviour. Rules, vectors, configuration and metadata are stored as versioned
artefacts, retaining separately usable association-based and embedding-based approaches. Training
experiments may reuse the same datasets without extracting them again.

Evaluation requires stored model artefacts and suitable versioned evaluation data. Before running
evaluation, the data scientist records the experiment protocol and acceptance criteria: the dataset
scope and split, list length, baseline, required metric thresholds for basket completion, and the
sample, relevance rubric and minimum acceptable result for the similar-product review. These
experiment settings are fixed before inspecting the candidate's evaluation results and retained
with the model and dataset versions. Publication checks the stored results against those criteria;
missing results or unmet criteria prevent publication. The thresholds are experiment parameters,
rather than evidence of online recommendation effectiveness.

Earlier purchases
train the models; later eligible multi-product orders form a chronological hold-out. Basket-completion
experiments expose part of each held-out basket as context and use its remaining products as targets.
Rule-only and hybrid recommendations are compared with a popularity baseline learned only from
training data, using precision and recall at a fixed list length and catalogue coverage across the
evaluated recommendation lists. The split, context/target procedure and list length are recorded.

Embedding-based similar products are assessed separately using a fixed sample of product queries
and a documented manual relevance rubric, with a category-based baseline evaluated on the same
queries. The review records the catalogue-text source, sample, relevance judgements and model version;
category-only Olist text limits the specificity of this initial semantic evaluation. Purchase
co-occurrence does not establish whether two products are suitable substitutes. The experiment
results support offline model comparison, without measuring recommendation-attributed sales.

Results are stored with the model and dataset versions. Publication requires acceptable stored
evaluation results under the documented acceptance criteria for the exact version being released;
it makes that version's rules, vectors, configuration and metadata available without rerunning
evaluation or automatically changing the active selection.

\paragraph{Refresh model experiments.}
When new completed purchases or catalogue updates become available, the data scientist manually runs
the warehouse refresh, extracts new versioned datasets and repeats the affected training and
evaluation experiments. New results are compared with the previous version before deciding whether
to publish and activate a replacement. Training-parameter trials may reuse an existing dataset.
This workflow reuses the warehouse, extraction, training and evaluation use cases; it does not
require interaction-event collection or automatic retraining.

\paragraph{Select and activate model/configuration.}
The data scientist chooses an available published model version or configuration. The system checks
that the selected artefacts are available, activates the selection for online serving and confirms
the active version/configuration. Subsequent recommendation requests use that selection, which the
administrator can inspect through \emph{View model status}. If activation fails, the current active
selection remains in use. Publication and activation may be performed consecutively by the data
scientist, who alone controls both operations.

% ---------------------------------------------------------------------
\clearpage
\section{Dataset Description and Project Integration}
\label{sec:dataset}

\subsection{Dataset overview}

The project will use the \emph{Brazilian E-Commerce Public Dataset by Olist},
available on Kaggle~\cite{olist2026}. Olist released anonymised commercial records
for approximately 100,000 orders placed between 2016 and 2018 across multiple
marketplaces in Brazil. The data describes customers, products, sellers, orders,
purchased items, payments, delivery information and customer reviews. It covers
multiple product categories, providing a historical basis for the project's
operational database, sales analytics and association-rule experiments.

The release consists of nine related CSV files. They form a relational data model:
orders connect customers to purchased items, and each item connects an order to a
product and a seller. Payments and reviews describe the same orders, while
geolocation and category translations provide additional reference information.
This connected structure allows the project to reconstruct purchase baskets and
prepare analytical datasets through the ETL pipeline.

\subsection{Files and fields}

The following tables give a simple description of all 52 fields across the nine
files. Attribute names follow the original CSV headers. The data types describe
how values will be represented after import; CSV files themselves do not declare
database types. String denotes text, Integer denotes a whole number, Decimal
denotes a numeric value with a fractional part, and Timestamp denotes a date and
time. Identifiers and postal-code prefixes will be stored as strings to preserve
their original form, including leading zeros. Missing values will be checked
during import rather than treated as zero.

\subsubsection{Customers}

\textbf{File:} \nolinkurl{olist_customers_dataset.csv}.

Customer identifiers and delivery locations. The same person can have different customer IDs across orders; the unique customer ID links repeat purchases.

{\tablestyle
\begin{longtable}{|L{5.7cm}|L{2.1cm}|L{\dimexpr\textwidth-7.8cm-6\tabcolsep-4\arrayrulewidth\relax}|}
  \caption{Fields in the customers file.}
  \label{tab:dataset-1}\\
  \hline
  \rowcolor{headbg}
  \textbf{Attribute} & \textbf{Data type} & \textbf{Description} \\[3pt]\hline
  \endfirsthead
  \hline
  \rowcolor{headbg}
  \textbf{Attribute} & \textbf{Data type} & \textbf{Description} \\[3pt]\hline
  \endhead
  \endfoot
  \endlastfoot
  \nolinkurl{customer_id} & String & Customer identifier used to link a customer record to an order. \\[3pt]\hline
  \nolinkurl{customer_unique_id} & String & Identifier for the same customer across different orders. \\[3pt]\hline
  \nolinkurl{customer_zip_code_prefix} & String & First five digits of the customer postal code. \\[3pt]\hline
  \nolinkurl{customer_city} & String & City of the customer delivery address. \\[3pt]\hline
  \nolinkurl{customer_state} & String & Brazilian state of the customer delivery address. \\[3pt]\hline
\end{longtable}
}

\subsubsection{Orders}

\textbf{File:} \nolinkurl{olist_orders_dataset.csv}.

Order headers, including the customer, order status and purchase and delivery dates.

{\tablestyle
\begin{longtable}{|L{5.7cm}|L{2.1cm}|L{\dimexpr\textwidth-7.8cm-6\tabcolsep-4\arrayrulewidth\relax}|}
  \caption{Fields in the orders file.}
  \label{tab:dataset-2}\\
  \hline
  \rowcolor{headbg}
  \textbf{Attribute} & \textbf{Data type} & \textbf{Description} \\[3pt]\hline
  \endfirsthead
  \hline
  \rowcolor{headbg}
  \textbf{Attribute} & \textbf{Data type} & \textbf{Description} \\[3pt]\hline
  \endhead
  \endfoot
  \endlastfoot
  \nolinkurl{order_id} & String & Identifier of the order. \\[3pt]\hline
  \nolinkurl{customer_id} & String & Links the order to its customer record. \\[3pt]\hline
  \nolinkurl{order_status} & String & Status of the order, such as delivered or canceled. \\[3pt]\hline
  \nolinkurl{order_purchase_timestamp} & Timestamp & Date and time when the order was placed. \\[3pt]\hline
  \nolinkurl{order_approved_at} & Timestamp & Date and time when the order payment was approved. \\[3pt]\hline
  \nolinkurl{order_delivered_carrier_date} & Timestamp & Date and time when the order was handed to the carrier. \\[3pt]\hline
  \nolinkurl{order_delivered_customer_date} & Timestamp & Date and time when the customer received the order. \\[3pt]\hline
  \nolinkurl{order_estimated_delivery_date} & Timestamp & Delivery date estimated at the time of purchase. \\[3pt]\hline
\end{longtable}
}

\subsubsection{Order items}

\textbf{File:} \nolinkurl{olist_order_items_dataset.csv}.

Items purchased within each order. An order can contain several items, possibly supplied by different sellers. An item is identified by the combination of order ID and item sequence.

{\tablestyle
\begin{longtable}{|L{5.7cm}|L{2.1cm}|L{\dimexpr\textwidth-7.8cm-6\tabcolsep-4\arrayrulewidth\relax}|}
  \caption{Fields in the order items file.}
  \label{tab:dataset-3}\\
  \hline
  \rowcolor{headbg}
  \textbf{Attribute} & \textbf{Data type} & \textbf{Description} \\[3pt]\hline
  \endfirsthead
  \hline
  \rowcolor{headbg}
  \textbf{Attribute} & \textbf{Data type} & \textbf{Description} \\[3pt]\hline
  \endhead
  \endfoot
  \endlastfoot
  \nolinkurl{order_id} & String & Links the item to its order. \\[3pt]\hline
  \nolinkurl{order_item_id} & Integer & Sequential item number within the order; not a quantity field. \\[3pt]\hline
  \nolinkurl{product_id} & String & Links the purchased item to its product. \\[3pt]\hline
  \nolinkurl{seller_id} & String & Links the purchased item to the seller fulfilling it. \\[3pt]\hline
  \nolinkurl{shipping_limit_date} & Timestamp & Deadline for the seller to hand the item to the logistics partner. \\[3pt]\hline
  \nolinkurl{price} & Decimal & Sale price of the item, excluding its freight charge. \\[3pt]\hline
  \nolinkurl{freight_value} & Decimal & Freight charge associated with the item. \\[3pt]\hline
\end{longtable}
}

\subsubsection{Order payments}

\textbf{File:} \nolinkurl{olist_order_payments_dataset.csv}.

Payment records for orders. An order may use more than one payment method, so several payment records can belong to the same order.

{\tablestyle
\begin{longtable}{|L{5.7cm}|L{2.1cm}|L{\dimexpr\textwidth-7.8cm-6\tabcolsep-4\arrayrulewidth\relax}|}
  \caption{Fields in the order payments file.}
  \label{tab:dataset-4}\\
  \hline
  \rowcolor{headbg}
  \textbf{Attribute} & \textbf{Data type} & \textbf{Description} \\[3pt]\hline
  \endfirsthead
  \hline
  \rowcolor{headbg}
  \textbf{Attribute} & \textbf{Data type} & \textbf{Description} \\[3pt]\hline
  \endhead
  \endfoot
  \endlastfoot
  \nolinkurl{order_id} & String & Links the payment record to its order. \\[3pt]\hline
  \nolinkurl{payment_sequential} & Integer & Sequence distinguishing payment records within an order. \\[3pt]\hline
  \nolinkurl{payment_type} & String & Payment method, such as credit card, voucher or boleto. \\[3pt]\hline
  \nolinkurl{payment_installments} & Integer & Number of installments selected for the payment. \\[3pt]\hline
  \nolinkurl{payment_value} & Decimal & Amount recorded for this payment. \\[3pt]\hline
\end{longtable}
}

\subsubsection{Products}

\textbf{File:} \nolinkurl{olist_products_dataset.csv}.

Product categories and physical attributes. The name and description length fields contain character counts, rather than the original product text. The spelling \texttt{lenght} is retained from the source.

{\tablestyle
\begin{longtable}{|L{5.7cm}|L{2.1cm}|L{\dimexpr\textwidth-7.8cm-6\tabcolsep-4\arrayrulewidth\relax}|}
  \caption{Fields in the products file.}
  \label{tab:dataset-5}\\
  \hline
  \rowcolor{headbg}
  \textbf{Attribute} & \textbf{Data type} & \textbf{Description} \\[3pt]\hline
  \endfirsthead
  \hline
  \rowcolor{headbg}
  \textbf{Attribute} & \textbf{Data type} & \textbf{Description} \\[3pt]\hline
  \endhead
  \endfoot
  \endlastfoot
  \nolinkurl{product_id} & String & Identifier of the product. \\[3pt]\hline
  \nolinkurl{product_category_name} & String & Product category name in Portuguese. \\[3pt]\hline
  \nolinkurl{product_name_lenght} & Integer & Number of characters in the product name; the name itself is absent. \\[3pt]\hline
  \nolinkurl{product_description_lenght} & Integer & Number of characters in the description; the description itself is absent. \\[3pt]\hline
  \nolinkurl{product_photos_qty} & Integer & Number of photographs associated with the product. \\[3pt]\hline
  \nolinkurl{product_weight_g} & Integer & Product weight in grams. \\[3pt]\hline
  \nolinkurl{product_length_cm} & Integer & Product length in centimetres. \\[3pt]\hline
  \nolinkurl{product_height_cm} & Integer & Product height in centimetres. \\[3pt]\hline
  \nolinkurl{product_width_cm} & Integer & Product width in centimetres. \\[3pt]\hline
\end{longtable}
}

\subsubsection{Sellers}

\textbf{File:} \nolinkurl{olist_sellers_dataset.csv}.

Identifiers and locations of sellers who fulfil the purchased items.

{\tablestyle
\begin{longtable}{|L{5.7cm}|L{2.1cm}|L{\dimexpr\textwidth-7.8cm-6\tabcolsep-4\arrayrulewidth\relax}|}
  \caption{Fields in the sellers file.}
  \label{tab:dataset-6}\\
  \hline
  \rowcolor{headbg}
  \textbf{Attribute} & \textbf{Data type} & \textbf{Description} \\[3pt]\hline
  \endfirsthead
  \hline
  \rowcolor{headbg}
  \textbf{Attribute} & \textbf{Data type} & \textbf{Description} \\[3pt]\hline
  \endhead
  \endfoot
  \endlastfoot
  \nolinkurl{seller_id} & String & Identifier of the seller. \\[3pt]\hline
  \nolinkurl{seller_zip_code_prefix} & String & First five digits of the seller postal code. \\[3pt]\hline
  \nolinkurl{seller_city} & String & City where the seller is located. \\[3pt]\hline
  \nolinkurl{seller_state} & String & Brazilian state where the seller is located. \\[3pt]\hline
\end{longtable}
}

\subsubsection{Order reviews}

\textbf{File:} \nolinkurl{olist_order_reviews_dataset.csv}.

Customer satisfaction survey responses about the purchase experience. Reviews are linked to orders, so an order review is not automatically a rating of each individual product.

{\tablestyle
\begin{longtable}{|L{5.7cm}|L{2.1cm}|L{\dimexpr\textwidth-7.8cm-6\tabcolsep-4\arrayrulewidth\relax}|}
  \caption{Fields in the order reviews file.}
  \label{tab:dataset-7}\\
  \hline
  \rowcolor{headbg}
  \textbf{Attribute} & \textbf{Data type} & \textbf{Description} \\[3pt]\hline
  \endfirsthead
  \hline
  \rowcolor{headbg}
  \textbf{Attribute} & \textbf{Data type} & \textbf{Description} \\[3pt]\hline
  \endhead
  \endfoot
  \endlastfoot
  \nolinkurl{review_id} & String & Identifier attached to the review record. \\[3pt]\hline
  \nolinkurl{order_id} & String & Links the review to its order. \\[3pt]\hline
  \nolinkurl{review_score} & Integer & Customer satisfaction score from 1 to 5. \\[3pt]\hline
  \nolinkurl{review_comment_title} & String & Optional title written by the customer. \\[3pt]\hline
  \nolinkurl{review_comment_message} & String & Optional review comment written by the customer. \\[3pt]\hline
  \nolinkurl{review_creation_date} & Timestamp & Date when the satisfaction survey was sent to the customer. \\[3pt]\hline
  \nolinkurl{review_answer_timestamp} & Timestamp & Date and time when the customer submitted the survey response. \\[3pt]\hline
\end{longtable}
}

\subsubsection{Geolocation}

\textbf{File:} \nolinkurl{olist_geolocation_dataset.csv}.

Geographic coordinates associated with Brazilian postal-code prefixes. Several location records can share the same prefix.

{\tablestyle
\begin{longtable}{|L{5.7cm}|L{2.1cm}|L{\dimexpr\textwidth-7.8cm-6\tabcolsep-4\arrayrulewidth\relax}|}
  \caption{Fields in the geolocation file.}
  \label{tab:dataset-8}\\
  \hline
  \rowcolor{headbg}
  \textbf{Attribute} & \textbf{Data type} & \textbf{Description} \\[3pt]\hline
  \endfirsthead
  \hline
  \rowcolor{headbg}
  \textbf{Attribute} & \textbf{Data type} & \textbf{Description} \\[3pt]\hline
  \endhead
  \endfoot
  \endlastfoot
  \nolinkurl{geolocation_zip_code_prefix} & String & First five digits of the postal code. \\[3pt]\hline
  \nolinkurl{geolocation_lat} & Decimal & Latitude of the location. \\[3pt]\hline
  \nolinkurl{geolocation_lng} & Decimal & Longitude of the location. \\[3pt]\hline
  \nolinkurl{geolocation_city} & String & City associated with the location. \\[3pt]\hline
  \nolinkurl{geolocation_state} & String & Brazilian state associated with the location. \\[3pt]\hline
\end{longtable}
}

\subsubsection{Category translations}

\textbf{File:} \nolinkurl{product_category_name_translation.csv}.

Portuguese-to-English translations of product category names. Categories without a matching translation will retain their original label.

{\tablestyle
\begin{longtable}{|L{5.7cm}|L{2.1cm}|L{\dimexpr\textwidth-7.8cm-6\tabcolsep-4\arrayrulewidth\relax}|}
  \caption{Fields in the category translations file.}
  \label{tab:dataset-9}\\
  \hline
  \rowcolor{headbg}
  \textbf{Attribute} & \textbf{Data type} & \textbf{Description} \\[3pt]\hline
  \endfirsthead
  \hline
  \rowcolor{headbg}
  \textbf{Attribute} & \textbf{Data type} & \textbf{Description} \\[3pt]\hline
  \endhead
  \endfoot
  \endlastfoot
  \nolinkurl{product_category_name} & String & Portuguese category name used in the products file. \\[3pt]\hline
  \nolinkurl{product_category_name_english} & String & Corresponding English category name. \\[3pt]\hline
\end{longtable}
}

\subsection{Project data model}
\label{sec:project-data-model}

Figure~\ref{fig:project-data-model} shows the adapted relational model for the
project. It separates a customer's persistent identity from the location snapshot
associated with an order, connects products to categories and sellers, and links
orders to their items, payments and reviews. The model retains purchase transactions and catalogue
information; it does not include an interaction-event entity. Current cart state and experiment
metadata are operational or analytical structures outside this identifier-focused diagram.

\begin{figure}[htbp]
  \centering
  \reportimage[width=\textwidth,height=0.75\textheight,keepaspectratio]{figures/data-model.png}{\textwidth}{0.69\textwidth}
  \caption{Adapted project data model and its entity relationships.}
  \label{fig:project-data-model}
\end{figure}
\clearpage

The following four tables describe the key attributes displayed in the diagram.
Attribute names are qualified by entity, for example \nolinkurl{Order.id}.
String is a text identifier, Int is an integer, and Bigint is a larger integer
used for generated record identifiers. PK denotes a primary key, FK a foreign
key, and UK a uniqueness constraint. Nullable means that the relationship can be
absent. The diagram focuses on identifiers and relationships; the descriptive,
transactional and timestamp fields listed in the CSV tables will also be mapped
to their corresponding project entities.

The source and project identifiers are distinct. Generated Int and Bigint keys
will identify records inside the project, while the original Olist identifiers
will be retained through source fields or mapping records. In particular,
\nolinkurl{customer_unique_id} will map to \nolinkurl{Customer.id}, and the
order-specific \nolinkurl{customer_id} will map to
\nolinkurl{CustomerSnapshot.id}. This preserves repeat-customer identity and the
customer location recorded for each historical order.

\paragraph{Customer identity and location.}

{\tablestyle
\begin{longtable}{|L{5.7cm}|L{2.1cm}|L{\dimexpr\textwidth-7.8cm-6\tabcolsep-4\arrayrulewidth\relax}|}
  \caption{Project model attributes: customer identity and location.}
  \label{tab:project-model-1}\\
  \hline
  \rowcolor{headbg}
  \textbf{Attribute} & \textbf{Data type} & \textbf{Description} \\[3pt]\hline
  \endfirsthead
  \hline
  \rowcolor{headbg}
  \textbf{Attribute} & \textbf{Data type} & \textbf{Description} \\[3pt]\hline
  \endhead
  \endfoot
  \endlastfoot
  \nolinkurl{Customer.id} & String & PK. Persistent customer identifier, mapped from the source customer unique ID. \\[3pt]\hline
  \nolinkurl{CustomerSnapshot.id} & String & PK. Identifies the customer location snapshot associated with a source order. \\[3pt]\hline
  \nolinkurl{CustomerSnapshot.customerId} & String & FK to Customer.id. Links the snapshot to the persistent customer identity. \\[3pt]\hline
  \nolinkurl{Geolocation.id} & Bigint & PK. Generated identifier for a geographic reference record. \\[3pt]\hline
  \nolinkurl{Geolocation.zipCodePrefix} & String & Postal-code prefix used to match customer and seller locations to geographic records. \\[3pt]\hline
\end{longtable}
}

\paragraph{Product catalogue and categories.}

{\tablestyle
\begin{longtable}{|L{5.7cm}|L{2.1cm}|L{\dimexpr\textwidth-7.8cm-6\tabcolsep-4\arrayrulewidth\relax}|}
  \caption{Project model attributes: product catalogue and categories.}
  \label{tab:project-model-2}\\
  \hline
  \rowcolor{headbg}
  \textbf{Attribute} & \textbf{Data type} & \textbf{Description} \\[3pt]\hline
  \endfirsthead
  \hline
  \rowcolor{headbg}
  \textbf{Attribute} & \textbf{Data type} & \textbf{Description} \\[3pt]\hline
  \endhead
  \endfoot
  \endlastfoot
  \nolinkurl{Category.id} & Int & PK. Generated identifier for a product category. \\[3pt]\hline
  \nolinkurl{CategoryTranslation.id} & Int & PK. Generated identifier for a category translation record. \\[3pt]\hline
  \nolinkurl{CategoryTranslation.categoryId} & Int & FK to Category.id and UK. Allows at most one translation record per category in this model. \\[3pt]\hline
  \nolinkurl{Product.id} & Int & PK. Generated product identifier, mapped to the original Olist product ID. \\[3pt]\hline
  \nolinkurl{Product.categoryId} & Int & FK to Category.id. Identifies the category classifying the product. \\[3pt]\hline
\end{longtable}
}

\paragraph{Orders, purchased items and sellers.}

{\tablestyle
\begin{longtable}{|L{5.7cm}|L{2.1cm}|L{\dimexpr\textwidth-7.8cm-6\tabcolsep-4\arrayrulewidth\relax}|}
  \caption{Project model attributes: orders, purchased items and sellers.}
  \label{tab:project-model-3}\\
  \hline
  \rowcolor{headbg}
  \textbf{Attribute} & \textbf{Data type} & \textbf{Description} \\[3pt]\hline
  \endfirsthead
  \hline
  \rowcolor{headbg}
  \textbf{Attribute} & \textbf{Data type} & \textbf{Description} \\[3pt]\hline
  \endhead
  \endfoot
  \endlastfoot
  \nolinkurl{Seller.id} & String & PK. Seller identifier retained from the source seller ID. \\[3pt]\hline
  \nolinkurl{Order.id} & Bigint & PK. Generated order identifier, mapped to the original Olist order ID. \\[3pt]\hline
  \nolinkurl{Order.customerId} & String & Nullable FK to Customer.id. Identifies the buyer when a persistent customer is available. \\[3pt]\hline
  \nolinkurl{Order.customerSnapshotId} & String & Nullable FK to CustomerSnapshot.id and UK. Associates an order with its own customer location snapshot. \\[3pt]\hline
  \nolinkurl{OrderItem.id} & Bigint & PK. Generated item-record identifier; the source order ID and item sequence are retained for mapping. \\[3pt]\hline
  \nolinkurl{OrderItem.orderId} & Bigint & FK to Order.id. Identifies the order containing the purchased item. \\[3pt]\hline
  \nolinkurl{OrderItem.productId} & Int & Nullable FK to Product.id. Identifies the purchased product when a linked catalogue record is available. \\[3pt]\hline
  \nolinkurl{OrderItem.sellerId} & String & Nullable FK to Seller.id. Identifies the seller fulfilling the item when available. \\[3pt]\hline
\end{longtable}
}

\paragraph{Payments and reviews.}

{\tablestyle
\begin{longtable}{|L{5.7cm}|L{2.1cm}|L{\dimexpr\textwidth-7.8cm-6\tabcolsep-4\arrayrulewidth\relax}|}
  \caption{Project model attributes: payments and reviews.}
  \label{tab:project-model-4}\\
  \hline
  \rowcolor{headbg}
  \textbf{Attribute} & \textbf{Data type} & \textbf{Description} \\[3pt]\hline
  \endfirsthead
  \hline
  \rowcolor{headbg}
  \textbf{Attribute} & \textbf{Data type} & \textbf{Description} \\[3pt]\hline
  \endhead
  \endfoot
  \endlastfoot
  \nolinkurl{OrderPayment.id} & Bigint & PK. Generated identifier for a payment record; its source order and payment sequence are retained. \\[3pt]\hline
  \nolinkurl{OrderPayment.orderId} & Bigint & FK to Order.id. Links a payment record to the order it pays for. \\[3pt]\hline
  \nolinkurl{OrderReview.id} & Bigint & PK. Generated review-record identifier, with the original review identifier retained for traceability. \\[3pt]\hline
  \nolinkurl{OrderReview.orderId} & Bigint & FK to Order.id. Links a satisfaction review to its purchase order. \\[3pt]\hline
\end{longtable}
}

Geolocation remains a reference entity in the supplied diagram; no direct
geolocation foreign key is shown. Its postal-code prefix will be used for the
location matching and consolidation described in the integration process below.

\subsection{Adapting the data model to the project}

We will adapt the Olist data model to the project's operational database and data
warehouse. The CSV files will first be imported into staging tables so that the
original records and identifiers remain available. Their entities will then be
mapped to the project's customers, products, categories, sellers, orders, order
items, payments and reviews. Geolocation and category translations will serve as
reference tables. Original identifiers will be retained for traceability, while
warehouse fact and dimension tables will use the keys and grain defined by the
project's analytical schema.

The main links will be preserved: \nolinkurl{customer_id} connects orders to
customer records; \nolinkurl{order_id} connects orders to items, payments and
reviews; and \nolinkurl{product_id} and \nolinkurl{seller_id} connect items to
products and sellers. \nolinkurl{customer_unique_id} will identify repeat buyers.
Category names will link products to translations, and postal-code prefixes will
link customer and seller locations to geography. Location records will be
consolidated by prefix before joining, and payment and review records will be
handled at their own grain to avoid multiplying order-item values in reports.

For Association Rule Mining, the project will group distinct product IDs within
eligible completed orders to construct purchase baskets. Purchase timestamps will
support chronological training and evaluation. Historical item prices and freight
charges will remain attached to their order items, rather than being overwritten
by current catalogue prices. The administration dashboard will use the resulting
warehouse records to calculate sales trends, order counts and product performance.
For the electronics storefront, a documented subset of relevant categories will be
selected while retaining the source order links and recording the scope of each
experimental dataset.

The source requires several adaptations to support the application. It does not provide original
product names or descriptions or inventory balances. Translated category text will support an
initial, limited semantic-recommendation baseline; richer catalogue text, if added for the project,
will be documented separately and linked to the same product IDs. Demonstration stock levels and
storefront display labels will be identified as project-created data.

New completed storefront purchases retain their order links, product identifiers, quantities,
purchase-time prices, timestamps and status. Together with catalogue updates, they can be loaded
into the warehouse alongside historical purchases during a manually initiated refresh. The data
scientist then creates a new dataset version and retrains and evaluates affected models before
publication or activation. Carts remain current shopping state; browsing, click, cart-action and
recommendation-impression histories are outside this project's data model. This approach fits the
dataset's relational structure into the project while distinguishing imported records, derived
information and data collected or created by the new system.

% ---------------------------------------------------------------------
% Insert before References. Upload the 15 referenced prototype-*.jpg images
% into figures/. Uses the existing reportimage helper and table packages.
% ---------------------------------------------------------------------
\clearpage
\section{Interactive Prototype and Interface Mockups}
\label{sec:prototype}

\emph{Signal Shop} demonstrates the buyer, administrator and data-scientist
interfaces for Section~\ref{sec:fr}. The figures use report-preview mode with a
curated electronics catalogue and sample orders. Recommendation pairings are
hardcoded, similar products use category matches, and model values are illustrative.

\subsection{Buyer interface and explained recommendations}
\label{sec:prototype-buyer}

The buyer flow is \emph{browse or search $\rightarrow$ product details
$\rightarrow$ recommendations $\rightarrow$ bag $\rightarrow$ checkout
$\rightarrow$ confirmation}. Table~\ref{tab:prototype-buyer} maps these screens to
Table~\ref{tab:fr-buyer}.

\begin{table}[H]
  \centering
  \caption{Buyer screens and related requirements.}
  \label{tab:prototype-buyer}
  \tablestyle
  \begin{tabularx}{\textwidth}{|L{3.1cm}|>{\RaggedRight\arraybackslash}X|L{2.7cm}|}
    \hline
    \rowcolor{headbg}
    \textbf{Screen} & \textbf{Information and actions} & \textbf{Requirements} \\\hline
    Storefront and catalogue & Featured products, category browsing and sorting
      (Figures~\ref{fig:prototype-storefront}--\ref{fig:prototype-catalogue}).
      & \req{FR-B01} \\\hline
    Search and details & Keyword results, product attributes, price and stock
      (Figures~\ref{fig:prototype-search-results}--\ref{fig:prototype-product-detail}).
      & \req{FR-B01} \\\hline
    Bought together & Four complementary items with explanations on the product
      page (Figure~\ref{fig:prototype-together}).
      & \req{FR-B04}, \req{FR-B06} \\\hline
    Similar products & Four alternatives below the selected product, retaining
      search context (Figure~\ref{fig:prototype-similar}).
      & \req{FR-B05}, \req{FR-B06} \\\hline
    Bag & Quantities, subtotal and suggestions that exclude selected products
      (Figure~\ref{fig:prototype-cart}).
      & \req{FR-B02}, \req{FR-B04}, \req{FR-B06} \\\hline
    Checkout and confirmation & Minimal contact details and a simulated order
      number (Figures~\ref{fig:prototype-checkout}--\ref{fig:prototype-confirmation}).
      & \req{FR-B02} \\\hline
  \end{tabularx}
\end{table}

\begin{samepage}
\subsubsection{Storefront and featured products}
\label{sec:prototype-storefront-view}

\textbf{Related requirements:} \req{FR-B01}.

Figure~\ref{fig:prototype-storefront} shows the main page and its featured
product list. Buyers can open an item or continue to the full collection.

\nopagebreak[4]
\begin{figure}[H]
  \centering
  \reportimage[width=\textwidth,height=0.62\textheight,keepaspectratio]{figures/prototype-storefront.jpg}{\textwidth}{0.56\textwidth}
  \caption{Main storefront with featured products and access to the collection.}
  \label{fig:prototype-storefront}
\end{figure}
\end{samepage}

\begin{samepage}
\subsubsection{Product catalogue}
\label{sec:prototype-catalogue-view}

\textbf{Related requirements:} \req{FR-B01}.

Figure~\ref{fig:prototype-catalogue} shows the collection page with product cards,
category filters and sorting. Each card provides a route to product details.

\nopagebreak[4]
\begin{figure}[H]
  \centering
  \reportimage[width=\textwidth,height=0.62\textheight,keepaspectratio]{figures/prototype-catalogue.jpg}{\textwidth}{0.56\textwidth}
  \caption{Product catalogue with category and sorting controls.}
  \label{fig:prototype-catalogue}
\end{figure}
\end{samepage}

\begin{samepage}
\subsubsection{Search results}
\label{sec:prototype-search-view}

\textbf{Related requirements:} \req{FR-B01}.

Searching for ``smartphone'' returns matching products
(Figure~\ref{fig:prototype-search-results}). The buyer selects Nova X1 to inspect
its details and recommendations.

\nopagebreak[4]
\begin{figure}[H]
  \centering
  \reportimage[width=\textwidth,height=0.62\textheight,keepaspectratio]{figures/prototype-search-results.jpg}{\textwidth}{0.56\textwidth}
  \caption{Keyword search results for smartphones.}
  \label{fig:prototype-search-results}
\end{figure}
\end{samepage}

\begin{samepage}
\subsubsection{Product details}
\label{sec:prototype-product-detail-view}

\textbf{Related requirements:} \req{FR-B01}, \req{FR-B02}.

Figure~\ref{fig:prototype-product-detail} shows Nova X1's description, tags, price
and stock. The search breadcrumb preserves the buyer's location, and
\emph{Add to bag} starts the purchase flow.

\nopagebreak[4]
\begin{figure}[H]
  \centering
  \reportimage[width=\textwidth,height=0.62\textheight,keepaspectratio]{figures/prototype-product-detail.jpg}{\textwidth}{0.56\textwidth}
  \caption{Selected smartphone with product information and an add-to-bag action.}
  \label{fig:prototype-product-detail}
\end{figure}
\end{samepage}

\begin{samepage}
\subsubsection{Frequently bought together}
\label{sec:prototype-together-view}

\textbf{Related requirements:} \req{FR-B04}, \req{FR-B06}.

On the product page, selecting \emph{Frequently bought together} reveals four
complementary items (Figure~\ref{fig:prototype-together}). The selected product
remains visible above the suggestions and their explanations.

\nopagebreak[4]
\begin{figure}[H]
  \centering
  \reportimage[width=\textwidth,height=0.62\textheight,keepaspectratio]{figures/prototype-frequently-bought-together.jpg}{\textwidth}{0.56\textwidth}
  \caption{Nova X1 product page with explained complementary suggestions.}
  \label{fig:prototype-together}
\end{figure}
\end{samepage}

\begin{samepage}
\subsubsection{Similar products}
\label{sec:prototype-similar-view}

\textbf{Related requirements:} \req{FR-B05}, \req{FR-B06}.

Selecting \emph{Similar products} displays four smartphone alternatives
(Figure~\ref{fig:prototype-similar}). Product details, the search breadcrumb,
prices and explanations make the suggestion context explicit.

\nopagebreak[4]
\begin{figure}[H]
  \centering
  \reportimage[width=\textwidth,height=0.62\textheight,keepaspectratio]{figures/prototype-similar-products.jpg}{\textwidth}{0.56\textwidth}
  \caption{Nova X1 product page with four similar smartphone alternatives.}
  \label{fig:prototype-similar}
\end{figure}
\end{samepage}

\begin{samepage}
\subsubsection{Shopping cart and basket recommendations}
\label{sec:prototype-cart-view}

\textbf{Related requirements:} \req{FR-B02}, \req{FR-B04}, \req{FR-B06}.

Figure~\ref{fig:prototype-cart} shows the selected phone and charger, quantity
controls and subtotal. Suggestions use the whole basket and exclude selected
products; the checkout action continues the flow.

\nopagebreak[4]
\begin{figure}[H]
  \centering
  \reportimage[width=\textwidth,height=0.62\textheight,keepaspectratio]{figures/prototype-cart.jpg}{\textwidth}{0.56\textwidth}
  \caption{Shopping bag with selected products, subtotal and basket suggestions.}
  \label{fig:prototype-cart}
\end{figure}
\end{samepage}

\begin{samepage}
\subsubsection{Checkout}
\label{sec:prototype-checkout-view}

\textbf{Related requirements:} \req{FR-B02}.

Figure~\ref{fig:prototype-checkout} collects the buyer's name and email and shows
the order total. The preview explicitly states that no payment is collected and
the sample order is not saved.

\nopagebreak[4]
\begin{figure}[H]
  \centering
  \reportimage[width=\textwidth,height=0.62\textheight,keepaspectratio]{figures/prototype-checkout.jpg}{\textwidth}{0.56\textwidth}
  \caption{Demo checkout with minimal contact information and order total.}
  \label{fig:prototype-checkout}
\end{figure}
\end{samepage}

\begin{samepage}
\subsubsection{Order confirmation}
\label{sec:prototype-confirmation-view}

\textbf{Related requirements:} \req{FR-B02}.

Figure~\ref{fig:prototype-confirmation} confirms the demonstration order with an
order number and total. The bag is cleared; preview confirmation does not change
stored orders or sales figures.

\nopagebreak[4]
\begin{figure}[H]
  \centering
  \reportimage[width=\textwidth,height=0.62\textheight,keepaspectratio]{figures/prototype-order-confirmation.jpg}{\textwidth}{0.56\textwidth}
  \caption{Simulated order confirmation after completing checkout.}
  \label{fig:prototype-confirmation}
\end{figure}
\end{samepage}

\clearpage
\subsection{Administration dashboard and order inspection}
\label{sec:prototype-admin}

Table~\ref{tab:prototype-admin} maps the administration views to
Table~\ref{tab:fr-admin}. Preview sales use consistent USD amounts and exclude
canceled orders. Catalogue forms are shown for interface review; saving and
withdrawal require live mode.

\begin{table}[H]
  \centering
  \caption{Administration screens and related requirements.}
  \label{tab:prototype-admin}
  \tablestyle
  \begin{tabularx}{\textwidth}{|L{3.1cm}|>{\RaggedRight\arraybackslash}X|L{2.7cm}|}
    \hline
    \rowcolor{headbg}
    \textbf{Screen} & \textbf{Information and actions} & \textbf{Requirements} \\\hline
    Catalogue and form & Product list, search, edit/withdraw actions and product
      data entry (Figures~\ref{fig:prototype-admin-catalogue}--\ref{fig:prototype-product-form}).
      & \req{FR-A01} \\\hline
    Sales overview & Revenue, orders, average order value, best sellers and
      trends with period selection (Figure~\ref{fig:prototype-sales}).
      & \req{FR-A03} \\\hline
    Orders & Search, status filtering, pagination and order selection
      (Figure~\ref{fig:prototype-orders}). & \req{FR-A02} \\\hline
    Order inspection & Customer, date, status, line items, purchase-time prices
      and totals (Figure~\ref{fig:prototype-order-inspection}).
      & \req{FR-A02}, \req{FR-B03} \\\hline
  \end{tabularx}
\end{table}

\begin{samepage}
\subsubsection{Catalogue management}
\label{sec:prototype-admin-catalogue-view}

\textbf{Related requirements:} \req{FR-A01}.

Figure~\ref{fig:prototype-admin-catalogue} shows the administrator's product
list with search, category filtering and edit/withdraw actions. Withdrawing a
product in live mode retains historical order items.

\nopagebreak[4]
\begin{figure}[H]
  \centering
  \reportimage[width=\textwidth,height=0.62\textheight,keepaspectratio]{figures/prototype-admin-catalogue.jpg}{\textwidth}{0.56\textwidth}
  \caption{Administrator product list with catalogue-management actions.}
  \label{fig:prototype-admin-catalogue}
\end{figure}
\end{samepage}

\begin{samepage}
\subsubsection{Product information form}
\label{sec:prototype-product-form-view}

\textbf{Related requirements:} \req{FR-A01}.

Figure~\ref{fig:prototype-product-form} shows sample data entered for a new
product. The form covers name, description, price, stock, category and tags;
editing uses the same fields.

\nopagebreak[4]
\begin{figure}[H]
  \centering
  \reportimage[width=\textwidth,height=0.62\textheight,keepaspectratio]{figures/prototype-product-form.jpg}{\textwidth}{0.56\textwidth}
  \caption{Product-entry form populated with demonstration values.}
  \label{fig:prototype-product-form}
\end{figure}
\end{samepage}

\begin{samepage}
\subsubsection{Sales overview}
\label{sec:prototype-sales-view}

\textbf{Related requirements:} \req{FR-A03}.

Figure~\ref{fig:prototype-sales} displays sales indicators and trends for the
selected seven-day or thirty-day period. These are sample-order aggregates;
warehouse-derived reporting remains future work.

\nopagebreak[4]
\begin{figure}[H]
  \centering
  \reportimage[width=\textwidth,height=0.62\textheight,keepaspectratio]{figures/prototype-sales.jpg}{\textwidth}{0.56\textwidth}
  \caption{Sales overview with period selection and sample indicators.}
  \label{fig:prototype-sales}
\end{figure}
\end{samepage}

\begin{samepage}
\subsubsection{Order browsing and selection}
\label{sec:prototype-orders-view}

\textbf{Related requirements:} \req{FR-A02}.

Figure~\ref{fig:prototype-orders} lets the administrator search, filter and
paginate orders. Selecting a record opens the inspection page in
Section~\ref{sec:prototype-order-inspection-view}.

\nopagebreak[4]
\begin{figure}[H]
  \centering
  \reportimage[width=\textwidth,height=0.62\textheight,keepaspectratio]{figures/prototype-orders.jpg}{\textwidth}{0.56\textwidth}
  \caption{Searchable order records linked to inspection pages.}
  \label{fig:prototype-orders}
\end{figure}
\end{samepage}

\begin{samepage}
\subsubsection{Order inspection}
\label{sec:prototype-order-inspection-view}

\textbf{Related requirements:} \req{FR-A02}, \req{FR-B03}.

Figure~\ref{fig:prototype-order-inspection} shows a sample order's customer,
status, purchase date and line items. In live mode, inspection reads stored
operational records and retains purchase-time names, quantities and prices.

\nopagebreak[4]
\begin{figure}[H]
  \centering
  \reportimage[width=\textwidth,height=0.62\textheight,keepaspectratio]{figures/prototype-order-inspection.jpg}{\textwidth}{0.56\textwidth}
  \caption{Sample order details with purchase-time quantities and prices.}
  \label{fig:prototype-order-inspection}
\end{figure}
\end{samepage}

\clearpage
\subsection{Data-scientist workspace}
\label{sec:prototype-scientist}

The retained overview illustrates version traceability. External scripts and
notebooks remain responsible for the eventual data preparation and training.

\begin{samepage}
\subsubsection{Dataset and model versions}
\label{sec:prototype-versions-view}

\textbf{Related requirements:} \req{FR-D02}, \req{FR-D06}, \req{FR-D10}.

Figure~\ref{fig:prototype-data-science} links sample model versions to dataset
references, preparation times and statuses. It visualises experiment metadata
without performing dataset extraction or retraining.

\nopagebreak[4]
\begin{figure}[H]
  \centering
  \reportimage[width=\textwidth,height=0.62\textheight,keepaspectratio]{figures/prototype-data-science.jpg}{\textwidth}{0.56\textwidth}
  \caption{Sample dataset and model versions in the data-scientist workspace.}
  \label{fig:prototype-data-science}
\end{figure}
\end{samepage}

\clearpage
\subsection{Prototype boundary and remaining implementation}
\label{sec:prototype-boundary}

The requirement links describe interface coverage, not completion of the
analytical requirements. Dedicated active-model status (\req{FR-A04}) and detailed
evaluation, publication and activation views (\req{FR-D07}--\req{FR-D09}) are omitted
from the mockups.

\begin{table}[H]
  \centering
  \caption{Prototype scope and remaining work.}
  \label{tab:prototype-boundary}
  \tablestyle
  \begin{tabularx}{\textwidth}{|L{2.6cm}|>{\RaggedRight\arraybackslash}X|>{\RaggedRight\arraybackslash}X|}
    \hline
    \rowcolor{headbg}
    \textbf{Area} & \textbf{Demonstrated} & \textbf{Remaining implementation} \\\hline
    Recommendations & Explained curated pairings, four alternatives and
      whole-basket context. & Learned rules, word embeddings and hybrid ranking
      (\req{FR-B04}, \req{FR-B05}, \req{FR-D03}--\req{FR-D05}). \\\hline
    Model overview & Sample dataset references, model versions and statuses.
      & Versioned datasets/artefacts, measured evaluation, publication, activation
      and role enforcement (\req{FR-D02}, \req{FR-D06}--\req{FR-D10}). \\\hline
    Purchases and catalogue & Buyer flow, order inspection and catalogue forms.
      & Preview checkout is not persisted; catalogue saving/withdrawal require
      live mode. \\\hline
    Reports and warehouse & Consistent sample sales; no warehouse-management
      mockup. & ETL, analytical structures and warehouse-derived reports
      (\req{FR-D01}, \req{FR-W01}--\req{FR-W06}). \\\hline
  \end{tabularx}
\end{table}

\clearpage
\phantomsection
\addcontentsline{toc}{section}{References}
\begin{thebibliography}{9}
  \bibitem{olist2026} Olist, \emph{Brazilian E-Commerce Public Dataset by Olist}.
    Kaggle. Available: \url{https://www.kaggle.com/datasets/olistbr/brazilian-ecommerce}.
    Accessed: 10 October 2026.
  \bibitem{agrawal1994} R.~Agrawal and R.~Srikant, ``Fast algorithms for mining association rules,''
    in \emph{Proc.\ 20th Int.\ Conf.\ Very Large Data Bases (VLDB)}, 1994, pp.~487--499.
  \bibitem{han2000} J.~Han, J.~Pei, and Y.~Yin, ``Mining frequent patterns without candidate
    generation,'' in \emph{Proc.\ ACM SIGMOD}, 2000, pp.~1--12.
  \bibitem{mikolov2013} T.~Mikolov, K.~Chen, G.~Corrado, and J.~Dean, ``Efficient estimation of word
    representations in vector space,'' \emph{arXiv:1301.3781}, 2013.
  \bibitem{bojanowski2017} P.~Bojanowski, E.~Grave, A.~Joulin, and T.~Mikolov, ``Enriching word vectors
    with subword information,'' \emph{Trans.\ Assoc.\ Comput.\ Linguistics}, vol.~5, pp.~135--146, 2017.
  \bibitem{ricci2015} F.~Ricci, L.~Rokach, and B.~Shapira, Eds., \emph{Recommender Systems Handbook},
    2nd~ed. Springer, 2015.
  \bibitem{cockburn2000} A.~Cockburn, \emph{Writing Effective Use Cases}. Addison-Wesley, 2000.
  \bibitem{kimball2013} R.~Kimball and M.~Ross, \emph{The Data Warehouse Toolkit: The Definitive Guide
    to Dimensional Modeling}, 3rd~ed. Wiley, 2013.
\end{thebibliography}

\end{document}
